\writeSectionFrame{\presentationTitle}{Das Null-Object-Pattern}
\begin{frame}{Lösung: Null Object Pattern}
    \begin{block}{Grundidee}
        Statt \texttt{null} wird ein spezielles Objekt verwendet, das nichts tut
        oder eine Null-Fall-Aktion ausführt.
    \end{block}
    \pause
    \begin{itemize}
        \item Kein null mehr $\rightarrow$ keine Nullprüfungen nötig
        \item Einheitliches Verhalten
        \item Einfach testbar und wartbar
    \end{itemize}
\end{frame}

\begin{frame}[fragile]{Basisinterface}
	\begin{exampleblock}{Weapon.java}
		\begin{lstlisting}
public interface Weapon {
    void attack();
}
	\end{lstlisting}
	\end{exampleblock}
\end{frame}

\begin{frame}[fragile]{Konkretes Weapon + Null Object}
	\begin{exampleblock}{Sword.java}
		\begin{lstlisting}
public class Sword implements Weapon {
    @Override
    public void attack() {
        System.out.println("Schwertangriff!");
    }
}
		\end{lstlisting}
	\end{exampleblock}
	
	\begin{exampleblock}{NullWeapon.java}
		\begin{lstlisting}
public class NullWeapon implements Weapon {
    @Override
    public void attack() {
        // tut nichts
    }
}
		\end{lstlisting}
	\end{exampleblock}
\end{frame}

\begin{frame}[fragile]{Player mit Null Object}
	\begin{exampleblock}{Player.java}
		\begin{lstlisting}
public class Player {
    private final Weapon weapon;

    public Player(Weapon weapon) {
        if (weapon == null) {
        	this.weapon = new NullWeapon();
        } else {
        	this.weapon = weapon;
        }
    }

    public void attack() {
        weapon.attack(); // nie null
    }
}
		\end{lstlisting}
	\end{exampleblock}
\end{frame}

\begin{frame}[fragile]{Null Object mit Alternativ-Aktion}
	\begin{exampleblock}{UnarmedWeapon.java}
		\begin{lstlisting}
public class UnarmedWeapon implements Weapon {
    @Override
    public void attack() {
        System.out.println("...");
    }
}
		\end{lstlisting}
	\end{exampleblock}
\end{frame}

